\title{Noesis: A Protocol for Trading and Fulfilling Machine Intelligence as a Commodity}

\author{Giacinto Paolo (GP) Saggese}
\thanks{University of Maryland, College Park, MD 20742, USA}

\date{\today}

% TODO(gp): Add the concept of max delay (e.g., to execute workload within N
% hours for batch execution)

% TODO(gp): People pay for token but they should pay for functionality.
% (virtual tokens, intelligence token)
% Introduce "virtual AI tokens" as tokens normalized by intelligence

% The fragmentation of the markets is not positive for consumers.
% TODO(gp): Users don't have a way to measure quality of service (e.g,
% model nerfing, low availability)

% TODO(gp): Add cache in the architecture if it's not there already

\begin{abstract}
Applications and autonomous agents increasingly consume large language model
(LLM) inference as a raw input, in the same way that factories consume
electricity. Today this input is purchased through static, per-model price
lists set unilaterally by a small number of providers, with no price
discovery, no verification that the delivered service matched what was paid
for, and no mechanism to route a request to only as much capability as it
actually requires.
%
We propose \Noesis, a protocol that treats machine
intelligence as a fungible commodity and organizes its exchange around two
cooperating components.
% TODO(gp): Market -> Exchange?
\NoesisMarket{} is a periodic call auction that matches buyers, who bid for a
task volume at a service class (capability, latency, and throughput bands)
with a minimum reliability, against sellers, who ask a price for delivering
it, clearing at a single uniform price per service class.
%
\NoesisServer{} is the fulfillment and monitoring layer: an API gateway that
executes the matched contracts against real model providers, meters whether
the delivered service class and reliability met the contracted terms,
and feeds the result back into the market as a reputation and pricing signal.
As a byproduct, \NoesisServer{} logs every request and estimates how much
capability each request actually needs, enabling difficulty-aware routing and
answer fusion that trade cost for accuracy in opposite directions.
We formalize the contract, bid, and ask notation; state the market-clearing
problem as a finite search over candidate clearing prices, each reducing to a
linear program over matched quantities, and establish that a clearing price
exists under standard monotonicity assumptions; and describe the
request-level fulfillment
statistics needed to attribute a shortfall to a specific seller rather than
to noise. We compare \Noesis{} against direct provider APIs, pass-through
gateways, and raw-compute spot markets, and lay out the open mechanism-design
questions and a phased implementation plan for a prototype.
\end{abstract}

\maketitle

\setcounter{tocdepth}{2}
\tableofcontents

% ##############################################################################
\section{Introduction}
\label{sec:introduction}

\Noesis{} is a protocol for exchanging machine intelligence, understood as LLM
inference capacity delivered at a stated service class and reliability,
between buyers who need it and sellers who supply it. It brings together
price discovery through an open auction, a contract that bundles quality
attributes rather than a single price-per-token, and a fulfillment layer that
measures and reports whether the contract was actually honored.

Today's LLM APIs price by the token, a unit that measures computation rather
than the functionality a buyer actually receives: two providers can charge
the same per-token rate while silently differing in the capability,
availability, or consistency of what that token buys. A buyer has no
standard way to detect that a model has been quietly downgraded (``model
nerfing'') or is intermittently unavailable, and each provider's
token-pricing scheme is incompatible with every other's, fragmenting the
market in a way that hurts, rather than helps, buyers trying to compare
offers. \Noesis{} addresses this by pricing and metering a \emph{(service class,
reliability)} bundle instead of a raw token count, so that what is
bought and verified is delivered functionality, not consumed computation.

% ==============================================================================
\subsection{Machine intelligence as a commodity}
\label{sec:intelligence_as_commodity}

Electricity, natural gas, and bandwidth all share a property that makes them
tradeable as commodities: any unit supplied by any seller is, for practical
purposes, interchangeable with any unit supplied by any other seller, once
quality attributes (voltage and frequency for electricity, calorific value for
gas, latency and jitter for bandwidth) are specified and metered. LLM inference
exhibits a version of the same property. A request answered at a stated
service class, with a stated reliability, is,
from the buyer's perspective, largely interchangeable across providers that can
meet those terms, even though the underlying models, weights, and hardware
differ completely.

Unlike electricity, however, the traded unit cannot be reduced to a single
scalar such as kilowatt-hours. Raw compute, measured in FLOPs or GPU-hours,
is observable but is not what a buyer wants to purchase: two providers can
spend the same compute on the same request and produce answers of very
different quality, and a buyer typically does not know, or care, how much
compute was spent to produce a given answer. The good that a buyer actually
wants to buy is better described as a bundle of two parts:

\begin{itemize}
  \item \textbf{Service class}, made of:
    \begin{itemize}
      \item \textbf{Capability band}: how good the answer is expected to be.
      \item \textbf{Latency band}: how long the answer takes.
      \item \textbf{Throughput band}: how fast output arrives.
    \end{itemize}
  \item \textbf{Reliability level}: how often the class is actually met.
\end{itemize}

For latency-insensitive batch workloads, the latency band can be \emph{batch},
e.g., each request answered within 24 hours.
Section~\ref{sec:bundling} returns to this point, which drives most of the
protocol's design.

% ==============================================================================
\subsection{The Noesis Components}
\label{sec:protocol_overview}

\Noesis{} is organized around two cooperating components.

\begin{itemize}
  \item \textbf{\NoesisMarket}: a periodic call auction that matches buyers
    and sellers into contracts. Every $T$ minutes, buyers submit bids and
    sellers submit asks, bucketed by service class; within each bucket, bids
    and asks are matched at a single uniform clearing price, subject to class
    and reliability compatibility, and asks of a higher class can fill bids
    for a lower one (Section~\ref{sec:noesis_market}).

  \item \textbf{\NoesisServer}: the fulfillment and monitoring layer. It is
    an LLM API gateway, modeled on multi-provider routers such as
    OpenRouter\footnote{\url{https://openrouter.ai}}, that executes the
    requests matched by a \NoesisMarket{} contract against the underlying
    providers, caches responses to identical or near-duplicate requests to
    avoid paying a provider twice for the same answer, logs every
    prompt/response pair, and measures whether the delivered service
    class and reliability met the contract terms. Violations are
    reported back to \NoesisMarket{} as a reputation and pricing signal
    (Section~\ref{sec:noesis_server}).
\end{itemize}

The two components close a loop that is missing from today's LLM API
market: a contract is metered and not merely sold, and the metering
result changes the terms under which the same seller can sell in the
future. \NoesisServer{} additionally reuses the logged traffic for two
purposes orthogonal to fulfillment:

\begin{itemize}
  \item \textbf{Difficulty-aware routing}: sends a request to the cheapest
    provider expected to satisfy it.
  \item \textbf{Answer fusion}: combines several providers' answers when
    higher reliability is worth the extra cost.
\end{itemize}

% ==============================================================================
\subsection{The Noesis architecture}
\label{sec:modularity}

Figure~\ref{fig:architecture} lays out the resulting architecture, with the
five pluggable components of Table~\ref{tab:components} shaded in teal:

% rendered_images:begin
% ```graphviz
% digraph NoesisArchitecture {
%   // ---------------------------------------------------------------
%   bgcolor="transparent";
%   pad="0.15";
%   splines=spline;
%   nodesep=0.30;
%   ranksep=0.50;
% 
%   node [shape=box,
%         style="rounded,filled",
%         fillcolor="#FFFFFF",
%         color="#9AA9B8",
%         penwidth=1.8,
%         fontname="Helvetica",
%         fontcolor="#243B53",
%         fontsize=12,
%         margin="0.22,0.14",
%         height=0.50];
% 
%   edge [color="#A3B1C0",
%         penwidth=1.3,
%         arrowhead=vee,
%         arrowsize=0.75,
%         fontname="Helvetica",
%         fontsize=10,
%         fontcolor="#7B8794"];
% 
%   // ---------------------------------------------------------------
%   // Bottom-to-top, so that the participants sit under the matching engine.
%   // As a consequence, cluster titles use labelloc="b" to be drawn at the top,
%   // and edges pointing leftward or downward are written reversed with dir=back.
%   rankdir=BT;
%   newrank=true;
% 
%   // actor : rose   |   process : blue   |   fulfillment monitoring : sage green   |   external : violet
% 
%   subgraph cluster_participants {
%     label     = "Market participants";
%     labelloc  = "t";
%     fontname  = "Helvetica-Bold";
%     fontsize  = 12.5;
%     fontcolor = "#3A4A5C";
%     style     = "rounded,filled";
%     fillcolor = "#F7F9FB";
%     color     = "#C7D0DA";
%     penwidth  = 1.0;
%     margin    = 16;
% 
%     Demand [label=<<b>1. Demand</b><br/><font point-size="10" color="#A34F6C">Applications / agents</font>>,
%             fillcolor="#F6E1E8", color="#D98CA8", fontcolor="#6B2A44"];
%     Supply [label=<<b>1. Supply</b><br/><font point-size="10" color="#946A2E">Model &amp; compute providers</font>>,
%             fillcolor="#FBEBD4", color="#D9A85F", fontcolor="#6B4517"];
%   }
% 
%   subgraph cluster_noesis_market {
%     label     = "NoesisMarket";
%     labeljust = "l";
%     labelloc  = "b";
%     fontname  = "Helvetica-Bold";
%     fontsize  = 12.5;
%     fontcolor = "#3A4A5C";
%     style     = "rounded,filled";
%     fillcolor = "#F7F9FB";
%     color     = "#C7D0DA";
%     penwidth  = 1.0;
%     margin    = 16;
% 
%     ReputationLoop [label=<<b>8. Reputation and feedback</b><br/><font point-size="10" color="#41719C">Pluggable</font>>,
%                     fillcolor="#D3E3F3", color="#7CA6CE", fontcolor="#1F4E79"];
%     MatchingEngine [label=<<b>2. Matching engine</b><br/><font point-size="10" color="#41719C">Pluggable</font>>,
%                     fillcolor="#D3E3F3", color="#7CA6CE", fontcolor="#1F4E79"];
%     PricingFeed    [label=<<b>3. Pricing dissemination</b><br/><font point-size="10" color="#41719C">Pluggable</font>>,
%                     fillcolor="#D3E3F3", color="#7CA6CE", fontcolor="#1F4E79"];
%     ReputationLoop -> MatchingEngine [label="Gates eligibility", style=dashed];
%     MatchingEngine -> PricingFeed [label="Cleared prices", style=dashed];
%   }
% 
%   subgraph cluster_noesis_server {
%     label     = "NoesisServer";
%     labelloc  = "b";
%     fontname  = "Helvetica-Bold";
%     fontsize  = 12.5;
%     fontcolor = "#3A4A5C";
%     style     = "rounded,filled";
%     fillcolor = "#F7F9FB";
%     color     = "#C7D0DA";
%     penwidth  = 1.0;
%     margin    = 16;
% 
%     CapabilityMeasure [label=<<b>6. Capability measurement</b><br/><font point-size="10" color="#41719C">Pluggable</font>>,
%                         fillcolor="#D3E3F3", color="#7CA6CE", fontcolor="#1F4E79"];
%     Gateway           [label=<<b>4. API gateway</b>>,
%                         fillcolor="#D3E3F3", color="#7CA6CE", fontcolor="#1F4E79"];
%     Metering          [label=<<b>7. Metering logic</b><br/><font point-size="10" color="#4F7A5A">Fulfillment monitoring</font>>,
%                         fillcolor="#DFEDE0", color="#8FB79A", fontcolor="#2E5A3D"];
%     Fusion            [label=<<b>5. Answer fusion</b><br/><font point-size="10" color="#41719C">Pluggable</font>>,
%                         fillcolor="#D3E3F3", color="#7CA6CE", fontcolor="#1F4E79"];
%     Metering -> Gateway [label="Logged traffic", dir=back];
%     Gateway -> Fusion [label="Multi-provider\nfan-out", style=dashed];
%     CapabilityMeasure -> Metering [label="Delivered\ncapability", style=dashed];
%   }
% 
%   ProviderAPIs        [label=<<b>Provider APIs</b><br/><font point-size="10" color="#6B52A8">Per-provider endpoints</font>>,
%                         fillcolor="#E9E1F7", color="#A88FD9", fontcolor="#45296B"];
%   IntelligenceMeasure [label=<<font point-size="10">Intelligence-measure providers</font>>,
%                         shape=note, fillcolor="#E9E1F7", color="#A88FD9", fontcolor="#45296B"];
% 
%   Demand -> MatchingEngine [label="Bids"];
%   Supply -> MatchingEngine [label="Asks"];
%   MatchingEngine -> Gateway [label="Contracts", penwidth=1.6, color="#8592A3"];
%   Gateway -> ProviderAPIs [label="Requests /\nresponses", dir=both];
%   Fusion -> ProviderAPIs [label="Fan-out", style=dashed];
%   CapabilityMeasure -> IntelligenceMeasure [label="Capability reference", style=dashed, dir=back];
%   ReputationLoop -> Metering [
%     label="Reputation &\npricing feedback",
%     style=dashed, penwidth=1.6,
%     color="#C0455B", fontcolor="#C0455B",
%     dir=back
%   ];
% 
%   // Rows, from bottom to top
%   { rank=same; Demand; Supply; }
%   { rank=same; ReputationLoop; MatchingEngine; PricingFeed; }
%   { rank=same; CapabilityMeasure; Metering; Gateway; Fusion; }
%   { rank=same; IntelligenceMeasure; ProviderAPIs; }
% 
%   // Legend: a graph label on top, so that all the entries stay in a single row.
%   labelloc="t";
%   labeljust="c";
%   // The outer invisible table adds a spacer row below the legend. The gray,
%   // square, dark-bordered box sets the legend apart from the diagram.
%   label=<
%     <table border="0" cellborder="0" cellspacing="0" cellpadding="0">
%       <tr><td>
%         <table border="2" color="#52606D" bgcolor="#E4E7EB" cellborder="0" cellspacing="8" cellpadding="5">
%           <tr>
%             <td><font face="Helvetica-Bold" point-size="12.5" color="#3A4A5C">Legend</font></td>
%             <td border="2" style="rounded" bgcolor="#F6E1E8" color="#D98CA8"><font face="Helvetica" point-size="10" color="#6B2A44">Demand</font></td>
%             <td border="2" style="rounded" bgcolor="#FBEBD4" color="#D9A85F"><font face="Helvetica" point-size="10" color="#6B4517">Supply</font></td>
%             <td border="2" style="rounded" bgcolor="#D3E3F3" color="#7CA6CE"><font face="Helvetica" point-size="10" color="#1F4E79">Noesis component (pluggable)</font></td>
%             <td border="2" style="rounded" bgcolor="#DFEDE0" color="#8FB79A"><font face="Helvetica" point-size="10" color="#2E5A3D">Fulfillment monitoring</font></td>
%             <td border="2" style="rounded" bgcolor="#E9E1F7" color="#A88FD9"><font face="Helvetica" point-size="10" color="#45296B">External reference</font></td>
%           </tr>
%         </table>
%       </td></tr>
%       <tr><td height="36"></td></tr>
%     </table>>;
% }
% ```
% label=fig:architecture
% caption=The \Noesis{} architecture, with its five pluggable components shaded in teal, and a color legend for demand/supply, Noesis components, fulfillment monitoring, and external references.
% rendered_images:end
% render_images:begin
\begin{figure}[H]
  \includegraphics[width=\linewidth]{01_introduction.tex.figs/01_introduction.1.png}
  \caption{The \Noesis{} architecture, with its five pluggable components shaded in teal, and a color legend for demand/supply, Noesis components, fulfillment monitoring, and external references.}
  \label{fig:architecture}
\end{figure}
% render_images:end

\begin{enumerate}
  \item \textbf{Demand and supply}: applications and agents submit bids, and
    model and compute providers submit asks, to \NoesisMarket{}'s matching
    engine.

  \item \textbf{Matching engine} (pluggable): clears compatible bids and asks
    (Definition~\ref{def:compatibility}) into contracts, which are dispatched
    to \NoesisServer{}'s API gateway. This matching is \NoesisMarket's only
    structural commitment. The periodic uniform-price call auction of
    Section~\ref{sec:noesis_market} is one instantiation of this component,
    and a continuous double auction or a posted-price mechanism
    (Section~\ref{sec:open_questions_market}) could be substituted without
    touching \NoesisServer.

  \item \textbf{Pricing dissemination} (pluggable): \NoesisMarket{} clears a
    price $p^*(q, t)$ per class and round (Problem~\ref{prob:clearing}), and
    that price is a useful signal beyond the bidders and sellers of the round
    that produced it. A pricing-dissemination interface, e.g., a real-time
    push feed or an on-chain event log
    (Section~\ref{sec:decentralized_extension}), publishes cleared prices to
    subscribers, buyers timing a future bid, other market makers, monitoring
    dashboards, without requiring them to poll \NoesisMarket{} or participate
    in a round.

  \item \textbf{API gateway}: proxies matched requests to the underlying
    providers through per-provider APIs.

  \item \textbf{Answer fusion} (pluggable): the gateway can fan a request out
    to several providers through this component when reliability outweighs
    cost. The aggregator $g$ of Equation~\eqref{eq:fusion_aggregator} is
    already defined as a pluggable function, majority vote, a verifier model,
    or a learned combiner (Section~\ref{sec:fusion}); self-consistency
    \cite{wang2022selfconsistency} and mixture-of-agents \cite{wang2024moa}
    are two published instantiations of the same interface
    (Section~\ref{sec:related_fusion}).

  \item \textbf{Capability measurement} (pluggable): how \NoesisServer{}
    estimates the capability tier $\hat c(x)$ actually delivered by a response
    (Section~\ref{sec:fulfillment_monitoring}). A verifier model, a benchmark
    probe, agreement with a stronger reference model
    (Definition~\ref{def:agreement_label}), or an external third-party
    evaluation service, e.g.,
    \href{https://artificialanalysis.ai}{artificialanalysis.ai}\footnote{\url{https://artificialanalysis.ai}}, are
    interchangeable implementations of the same interface: a mapping from a
    prompt/response pair to a tier in $\calC$.

  \item \textbf{Metering logic}: the part of the gateway that checks the
    responses against the contract terms and reports fulfillment back to
    \NoesisMarket{}.

  \item \textbf{Reputation and feedback loop} (pluggable): turns the
    fulfillment reports of the metering logic into a score that gates future
    eligibility in the matching engine (Section~\ref{sec:reputation}). The
    exponential-decay update of Equation~\eqref{eq:reputation_update} is one
    instantiation of this component; any function that maps a seller's
    compliance history to such a score satisfies the same interface.
\end{enumerate}

Table~\ref{tab:components} summarizes the five pluggable components,
their responsibility, and the instantiation used elsewhere in this paper versus a
pluggable alternative. This independence is also what lets the matching
engine and the gateway-plus-logging layer be built and demonstrated as
separate, independently testable pieces of a prototype, precisely because
each is a separately pluggable component rather than part of one monolith.

\begin{samepage}
  \begin{table}[H]
    \centering
    \scriptsize
    \begin{tabular}{>{\raggedright\arraybackslash}p{2.4cm}|p{3.3cm}|c|p{2.9cm}|p{2.9cm}|}
      \textbf{Component} & \textbf{What it does} & \textbf{Sec.} & \textbf{This paper implementation} & \textbf{Alternative implementation} \\
      \hline
      \textbf{2. Matching engine} & Matches compatible bids and asks into contracts within a class bucket (Problem~\ref{prob:clearing}) & \ref{sec:noesis_market} & Periodic uniform-price call auction & Continuous double auction \\
      \hline
      \textbf{3. Pricing dissemination} & Publishes each round's cleared prices to subscribers outside the round that produced them & \ref{sec:modularity} & A real-time push feed & An on-chain event log (Section~\ref{sec:decentralized_extension}) \\
      \hline
      \textbf{5. Answer fusion} & Combines several providers' answers into one when reliability outweighs cost & \ref{sec:fusion} & Majority vote or a verifier model & A learned combiner, e.g., mixture-of-agents \cite{wang2024moa} \\
      \hline
      \textbf{6. Capability measurement} & Estimates the capability tier $\hat c(x)$ actually delivered by a response & \ref{sec:fulfillment_monitoring} & A verifier model, benchmark probe, or reference-model agreement (Definition~\ref{def:agreement_label}) & A hosted evaluation provider, e.g., \href{https://artificialanalysis.ai}{artificialanalysis.ai} \\
      \hline
      \textbf{8. Reputation and feedback loop} & Aggregates a seller's fulfillment history into a score that gates future eligibility & \ref{sec:reputation} & Exponential-decay update (Equation~\eqref{eq:reputation_update}) & A Bayesian trust score \\
      \hline
    \end{tabular}
    \caption{The five pluggable components behind \NoesisMarket{} and
      \NoesisServer, their responsibility, and the instantiation used
      elsewhere in this paper versus a pluggable alternative.}
    \label{tab:components}
  \end{table}
\end{samepage}

% ==============================================================================
\subsection{Contributions of This Paper}
\label{sec:contributions}

The innovative features of \Noesis{} are:

\begin{enumerate}
  \item \textbf{Modular, pluggable components}: \NoesisMarket{} and
    \NoesisServer{} are thin coordination layers around five independently
    replaceable components, each defined by a narrow interface rather than a
    fixed implementation (Section~\ref{sec:modularity},
    Table~\ref{tab:components}):
    \begin{itemize}
      \item \textbf{Matching engine}
      \item \textbf{Capability measurement}
      \item \textbf{Reputation and feedback}
      \item \textbf{Answer fusion}
      \item \textbf{Pricing dissemination}
    \end{itemize}

  \item \textbf{Price discovery through auction, not price lists}:
    \NoesisMarket{} clears buyers and sellers through a periodic call
    auction instead of relying on prices unilaterally posted by a small
    number of providers.

  \item \textbf{A bundled, verifiable good}: the traded unit is a
    (service class, reliability) bundle rather than a raw compute
    quantity, so the market quotes and clears on ``how good'', not only
    ``how much''.

  \item \textbf{Verified fulfillment}: \NoesisServer{} meters every
    request executed under a contract and computes a measured reliability
    from per-request capability, latency, and throughput checks, so a
    contract's terms are checked against evidence rather
    than taken on faith.

  \item \textbf{Reputation and pricing feedback}: a seller whose measured
    performance falls short of its ask is flagged, and the resulting
    reputation signal feeds back into eligibility and pricing in future
    auction rounds, in the spirit of non-performance penalties in capacity
    markets.

  \item \textbf{Difficulty-aware routing and answer fusion}: two
    complementary modes of the same request-shaping capability:
    \begin{itemize}
      \item \textbf{Difficulty-aware routing}: \NoesisServer{} estimates how
        much capability a request actually needs and routes accordingly,
        spending intelligence only where it is needed.
      \item \textbf{Answer fusion}: for requests where reliability matters
        more than cost, \NoesisServer{} can instead query several providers
        and combine their answers, trading cost for accuracy in the direction
        opposite to routing.
    \end{itemize}

  \item \textbf{Provider-agnostic liquidity pooling}: many sellers compete
    for the same demand behind a single interface, instead of buyers being
    locked into a single provider's price list.

  \item \textbf{Tiered commodity pricing}: each service class
    clears at its own price in each auction round, and higher classes can
    fill lower ones, analogous to peak and off-peak
    pricing in electricity markets.
\end{enumerate}

% ==============================================================================
\subsection{Comparison with alternative solutions}
\label{sec:comparison_table}

Table~\ref{tab:comparison} compares \Noesis{} against three alternative ways
of obtaining LLM inference:

\begin{itemize}
  \item \textbf{Direct API}: calling a provider's API directly.
  \item \textbf{Gateway}: using a pass-through multi-provider gateway (e.g.,
    OpenRouter-style routers).
  \item \textbf{Spot compute}: buying raw compute on a cloud spot market
    (e.g., spot GPU instances).
\end{itemize}

The comparison is discussed in detail in Sections~\ref{sec:market_design}
through \ref{sec:noesis_server}.

\begin{samepage}
  \begin{table}[H]
    \centering
    \begin{tabular}{l|c|c|c|c|}
                                                & \emph{Direct API} & \emph{Gateway} & \emph{Spot compute} & \emph{\Noesis} \\
      \hline
      Price discovery via open auction       & \nbad             & \nbad          & \ngood               & \ngood         \\
      \hline
      Bundled quality guarantee              & \npartial         & \npartial      & \nbad                & \ngood         \\
      \hline
      Verified fulfillment / SLA monitoring  & \nbad             & \nbad          & \npartial             & \ngood         \\
      \hline
      Reputation-based accountability        & \nbad             & \nbad          & \npartial             & \ngood         \\
      \hline
      Difficulty-aware cost routing          & \nbad             & \npartial      & \nbad                 & \ngood         \\
      \hline
      Multi-provider liquidity pooling       & \nbad             & \ngood         & \npartial             & \ngood         \\
      \hline
      Reusable evaluation/training dataset   & \npartial         & \npartial      & \nbad                 & \ngood         \\
      \hline
      Single-API ease of integration         & \ngood            & \ngood         & \nbad                 & \ngood         \\
      \hline
    \end{tabular}
    \caption{Comparison of \Noesis{} against alternative ways of obtaining
      LLM inference.}
    \label{tab:comparison}
  \end{table}
\end{samepage}

\begin{itemize}
  \item \textbf{Cloud spot markets} already apply auction-based price
    discovery, but to raw compute rather than to a service-class/reliability
    bundle, so they score poorly on the bundled-quality and
    difficulty-routing rows: a spot GPU market has no notion of ``how
    good'' the resulting answer is.
  \item \textbf{Gateways} pool liquidity across providers and are easy to
    integrate, but set prices unilaterally rather than clearing an
    auction, and typically do not meter fulfillment against a contract.
  \item \textbf{\Noesis{}} is designed to combine the liquidity pooling and
    ease-of-integration of a gateway with the price discovery of an
    auction market and the verifiability of a metered service contract.
\end{itemize}

Section~\ref{sec:state_of_the_art} surveys the current industry landscape,
startups and projects that address pieces of this same problem, in more
detail.
